\section{Wer mehr wissen will}
\subsection*{Wer mehr wissen will}


\begin{frame}{Mobile Geräte}
	\begin{itemize}
		\item K9 Mail
		\item APG
		\item Anleitung: %TODO url raussuchen
	\end{itemize}
\end{frame}

\begin{frame}{Facebook und PGP}
	\begin{itemize}
		%TODO recherschieren, links
		\item verschlüsselte Mails für
			\begin{itemize}
				\item Passwort zurücksetzen
				\item Notifications
			\end{itemize}
		\item Nachteil: Alle IDs / E-Mail Addressen bekannt
		\item Lösung: Nutze spezielle Addresse nur für FB
	\end{itemize}
	\note[item]{Nur eine PR Masche}
\end{frame}

\begin{frame}{Outlook und Windows Mail}
	\begin{itemize}
		% TODO Links, informieren
		\item Outlook hat Plugin
		\item Windows Live Mail nicht damit funktionabel
	\end{itemize}
\end{frame}

\begin{frame}{Links und weitere Veranstaltungen}
	\begin{itemize}
		\item \url{http://cryptoparty-hd.de} \\[10pt]
		\item \textbf{Anonymität} am Mo 29.06. INF 368, Raum 432
	\end{itemize}
\end{frame}

\begin{frame}{Bilder und Lizenzen}
	\smaller
	\begin{itemize}
		\item Von \url{https://en.wikipedia.org/wiki/File:Asymmetric_cryptography_-_step_2.svg}
			\begin{center} \includegraphics[height=0.7cm]{Bob} \hspace{10pt} \includegraphics[height=0.7cm]{Alice} \end{center}

		\item Vom ``Tango Icon Theme'' --- \url{http://tango.freedesktop.org} 
			\begin{center} \includegraphics[height=0.7cm]{Document}\hspace{10pt}  \includegraphics[height=0.7cm]{mail_open}\hspace{10pt}  \includegraphics[height=0.7cm]{mail_closed} \end{center}
		\item Von der HGS Mathcomp --- \url{http://mathcomp.org}
			\begin{center} \includegraphics[height=0.7cm]{logoHGS} \end{center}
		\item Alle anderen Bilder sind eigene Werke oder Derivate
		\item Letztere werden unter CC by-sa 4.0 veröffentlicht
	
	\vspace \fill
	\begin{center}
		\small
		\includegraphics[scale=0.5]{cc-by-sa}\\
		\scalebox{0.7}{This work is licensed under a Creative Commons}\\[-4pt]
		\scalebox{0.7}{Attribution-ShareAlike 4.0 International Licence.}
	\end{center}
	\end{itemize}
\end{frame}
